\documentclass[11pt]{article}
\usepackage[utf8]{inputenc}
\usepackage[spanish]{babel}
\usepackage{amsmath,amsthm,amsfonts,amssymb,amscd}
\usepackage{multirow,booktabs}
\usepackage[table]{xcolor}
\usepackage{fullpage}
\usepackage{lastpage}
\usepackage{enumitem}
\usepackage{fancyhdr}
\usepackage{mathrsfs}
\usepackage{wrapfig}
\usepackage{setspace}
\usepackage{calc}
\usepackage{multicol}
\usepackage[retainorgcmds]{IEEEtrantools}
\usepackage[margin=3cm]{geometry}
\usepackage{empheq}
\usepackage{framed}
\usepackage[most]{tcolorbox}
\usepackage{soul}
\usepackage{array}
\usepackage{ragged2e}
\setlength{\parindent}{0.0in}
\setlength{\parskip}{0.05in}
\colorlet{shadecolor}{orange!15}
\parindent 0in
\parskip 12pt
\geometry{margin=1in, headsep=0.25in}

\begin{document}
\setcounter{section}{0}

\thispagestyle{empty}

\begin{center}
{\LARGE \bf El Mayor Riesgo de Mi Tesis}\\[4pt]
{\large Tesis I}\\[2pt]
Roberto Treviño Cervantes
\end{center}

\section{Reflexión}

\textit{El mayor riesgo de mi tesis es} la validez del uso de una red neuronal convolucional (CNN) entrenada con el conjunto WM-811K como función de costo sustituta del \textit{yield} dentro del lazo de optimización de trayectorias. WM-811K es un repositorio de mapas de obleas (\textit{wafer bin maps}) provenientes de líneas de fabricación reales, sin trazabilidad sobre los perfiles de movimiento que originaron cada patrón de defecto. Por lo tanto, la red aprende a clasificar morfologías agregadas de defecto, pero no relaciona explícitamente una distribución de error de seguimiento con su mapa de \textit{yield} resultante. Si esta correspondencia no se sostiene, la función de costo aprendida puede inducir al PSO a optimizar trayectorias hacia regiones del espacio paramétrico que mejoran una métrica simulada sin reflejo en la calidad real del patrón impreso.

La mitigación se construye en tres capas. Primero, se acota explícitamente el alcance de la contribución a la \textit{optimización de trayectorias consciente del \textit{yield} dentro de la simulación}, sin mención del \textit{yield} de una fábrica real. Segundo, se diseña un protocolo de validación interno: generar mapas sintéticos de defecto a partir de perturbaciones controladas de trayectoria sobre el gemelo digital de Al-Rawashdeh y verificar que la CNN reproduzca las clases morfológicas correspondientes de WM-811K. Tercero, los resultados se reportan como una correspondencia \emph{cualitativa} entre regímenes cinemáticos y firmas de defecto, de modo que la conclusión metodológica permanezca sólida aun si la calibración cuantitativa al \textit{yield} de una fábrica requiere trabajo experimental futuro.

\end{document}
